\documentclass{article}
\usepackage[T1]{fontenc}
\usepackage{lmodern}
\usepackage{graphicx}
\usepackage{amsmath,amssymb}
\usepackage[a4paper,margin=2cm]{geometry}
\usepackage[absolute,overlay]{textpos}
\setlength{\TPHorizModule}{1mm}
\setlength{\TPVertModule}{1mm}
\usepackage[indonesian]{babel}
\usepackage{hyperref}
\setlength{\emergencystretch}{2em}
% unit: Halaman 12

\begin{document}

% === Halaman 12 (llm) ===

\section*{Medan Berhingga $F_p$}

\begin{itemize}
    \item Kelas medan berhingga yang penting adalah $F_p$, dengan $p$ adalah bilangan prima.
    \item $F_p$ adalah adalah medan berhingga dengan himpunan bilangan bulat $\{0, 1, 2, \dots, p-1\}$, $p$ bilangan prima, dan dua operasi yang didefinisikan sbb:
\end{itemize}

\begin{enumerate}
    \item \textbf{Penjumlahan}, dilakukan dalam modulus $p$
    \[
    \text{Jika } a, b \in F_p, \text{ maka } a + b = r, \text{ yang dalam hal ini}
    \]
    \[
    r = (a + b) \pmod p, \quad 0 \le r \le p - 1
    \]
    \item \textbf{Perkalian}, dilakukan dalam modulus $p$
    \[
    \text{Jika } a, b \in F_p, \text{ maka } a \cdot b = s, \text{ yang dalam hal ini}
    \]
    \[
    s = (a \cdot b) \pmod p, \quad 0 \le s \le p - 1
    \]
\end{enumerate}

\vfill
\begin{center}
    Rinaldi Munir/II4021 Kriptografi
\end{center}

\end{document}
